\subsection{Obstacles with bounded metric Laplacian}
\label{sec:metric-laplacian-obstacles}
The results above show that the Monge--Amp\`ere measure of an envelope is carried by the contact set, but they do not identify the measure there. For obstacles with bounded Laplacian we now compute the measure on the contact set from the second derivatives of the obstacle. Recall from \cref{sec:global-trace-equation} that $L=\mathrm d^*\mathrm d$ is the Laplace--Beltrami operator of the Riemannian metric defined by the fixed K\"ahler form $\omega$; with our sign convention it is nonnegative, and for the Euclidean metric it equals $-\Delta$. We call $L$ the \emph{metric Laplacian}. A function $f$ on an open set $U\subset X$ has \emph{(locally) bounded metric Laplacian} if the distribution $Lf$ on $U$ is a (locally) bounded measurable function.
\index{Laplacian!metric}\index{Laplacian!bounded metric Laplacian}

By the trace identity $\Lambda\ddc f=-Lf/(4\pi)$ of \cref{lem:kahler-trace-star}, where $\Lambda$ is the trace with respect to $\omega$, a bound for $Lf$ is a bound for the trace of the complex Hessian of $f$. For a $\theta$-psh function $u$ the current $\theta+\ddc u$ is positive, and positivity turns this trace bound into bounds for all complex Hessian coefficients. A general obstacle has no such positivity, and its complex Hessian need not be bounded. Nevertheless it has square-integrable weak second derivatives (\cref{lem:metric-laplacian-local-weak-hessian}), and on the set where it agrees with a potential that also has bounded metric Laplacian, these coincide almost everywhere with the second derivatives of the potential. This gives the contact formula of \cref{prop:metric-laplacian-regular-contact-density}. The $L^p$ estimate of \cref{prop:compact-metric-laplacian-lp} shows moreover that such an obstacle has weak second derivatives in every $L^p$ with $p<\infty$, so that wedge products of the corresponding weak forms are integrable on all of $X$ (\cref{cor:metric-laplacian-barrier-determinants}).

The first lemma is the variable-coefficient analogue of \cref{lem:bounded-laplacian-weak-hessian}. We use the Sobolev spaces $W^{k,p}$ of \cref{sec:hessian-singular-integrals}. A function on an open set $U\subset X$ belongs to $W^{2,2}_{\loc}(U)$ if it belongs to $W^{2,2}_{\loc}$ in every holomorphic chart of $U$. For an open set $U'\subset X$, the norm $\|f\|_{W^{k,p}(U')}$ is computed with the charts and the partition of unity that define the norm of $W^{k,p}(X)$: it is the sum, over the charts, of the $W^{k,p}$ norms of the products of $f$ with the corresponding cutoff functions, taken on the images of $U'$ in the charts.

\begin{lemma}\label{lem:metric-laplacian-local-weak-hessian}
Let $U\subset X$ be open, and let $f\in L^2_{\loc}(U)$ be such that its distributional metric Laplacian $Lf$ belongs to $L^2_{\loc}(U)$. Then $f\in W^{2,2}_{\loc}(U)$. More precisely, for open sets $U'\Subset V\Subset U$ there is a constant $C$, depending only on $U'$, $V$ and $\omega$, such that
\[
 \|f\|_{W^{2,2}(U')}\leq C\bigl(\|f\|_{L^2(V)}+\|Lf\|_{L^2(V)}\bigr).
\]
In particular, every continuous function on $U$ with locally bounded metric Laplacian belongs to $W^{2,2}_{\loc}(U)$.
\end{lemma}
\begin{proof}
We use the Sobolev spaces $H^r$ of \cref{sec:global-trace-equation}. Fix $c>0$. By \cref{lem:compact-elliptic-foundations}(3), applied to functions, on which $\mathcal L=L$, the operator $P=L+c$ is a bijection of the space of distributions on $X$ onto itself, and its inverse maps $H^{-1}$ boundedly into $H^1$ and $H^0=L^2$ boundedly into $H^2$. Since no derivative of $f$ is known in advance, we gain regularity in two steps, using two cutoff functions.

Choose $\chi_0\in C_c^\infty(V)$ equal to one near $\overline{U'}$, and $\chi_1\in C_c^\infty(V)$ supported in the interior of $\{\chi_0=1\}$ and also equal to one near $\overline{U'}$. Extend $\chi_0f$ by zero to a function on $X$. The product rule for distributions gives
\[
 P(\chi_0f)=\chi_0(Lf+cf)+[L,\chi_0]f.
\]
The commutator $[L,\chi_0]$ is a first-order differential operator with smooth coefficients supported in $V$. Pairing $[L,\chi_0]f$ with a smooth function $\phi$ and moving the derivative onto $\phi$, we bound the pairing by $C\|f\|_{L^2(V)}\|\phi\|_{H^1}$. Hence $[L,\chi_0]f\in H^{-1}$, and the right side of the last identity has $H^{-1}$ norm at most $C(\|Lf\|_{L^2(V)}+\|f\|_{L^2(V)})$. Applying $P^{-1}$, we obtain $\chi_0f\in H^1$ with the same bound.

In particular $f$ has square-integrable first derivatives on the interior of $\{\chi_0=1\}$, which contains the support of $\chi_1$. Therefore, in the identity
\[
 P(\chi_1f)=\chi_1(Lf+cf)+[L,\chi_1]f,
\]
the commutator term involves only $f$ and its first derivatives on the support of $\mathrm d\chi_1$, and the right side lies in $L^2$, with norm at most $C(\|Lf\|_{L^2(V)}+\|f\|_{L^2(V)})$. Since $P^{-1}$ is bounded from $L^2$ to $H^2$, it follows that $\chi_1f\in H^2$ with this bound. In each chart, $H^2$ coincides with $W^{2,2}$ by Plancherel's theorem (\cref{thm:plancherel}), as observed in Step 1 of the proof of \cref{lem:compact-elliptic-foundations}. Since $\chi_1=1$ near $\overline{U'}$, this proves the estimate, and covering $U$ by such sets $U'$ proves that $f\in W^{2,2}_{\loc}(U)$.
\end{proof}
\index{Laplacian!metric weak Hessian}

For the contact formula we work in a coordinate ball. Recall that $D(A_1,\ldots,A_n)$ denotes the mixed discriminant of Hermitian $n\times n$ matrices (\cref{sec:positive-currents}), that $\mathrm dV_{2n}$ is Lebesgue measure in the coordinates, and that $(u_{a\bar b})$, with $u_{a\bar b}=\partial_{z_a}\partial_{\bar z_b}u$, is the complex Hessian matrix of $u$ (\cref{sec:complex-exterior-calculus}).

\begin{proposition}\label{prop:metric-laplacian-regular-contact-density}
Let $U\subset X$ be open. For $1\leq j\leq n$, let $\theta_j$ be a closed smooth real $(1,1)$-form on $U$, let $u_j$ be a continuous $\theta_j$-psh function on $U$, and let $f_j$ be a continuous real function on $U$. Suppose that $u_j$ and $f_j$ have locally bounded metric Laplacian, and put
\[
 E=\bigcap_{j=1}^n\{u_j=f_j\}.
\]
Let $W\Subset U$ be a coordinate ball, choose smooth real functions $\rho_j$ on $W$ with $\ddc\rho_j=\theta_j$, and let
\[
 A_j=\bigl((\rho_j+u_j)_{a\bar b}\bigr)_{a,b},
 \qquad
 B_j=\bigl((\rho_j+f_j)_{a\bar b}\bigr)_{a,b}
\]
be the weak complex Hessian matrices, defined almost everywhere on $W$. Then each $A_j$ is positive semidefinite almost everywhere and locally bounded, and on $W$
\[
 \begin{aligned}
 \bigwedge_{j=1}^n(\theta_j+\ddc u_j)
       &=\frac{n!}{\pi^n}D(A_1,\ldots,A_n)\,\mathrm dV_{2n},\\
 \mathbf1_E\bigwedge_{j=1}^n(\theta_j+\ddc u_j)
       &=\frac{n!}{\pi^n}\mathbf1_E D(B_1,\ldots,B_n)\,\mathrm dV_{2n}.
 \end{aligned}
\]
In particular the density $\mathbf1_ED(B_1,\ldots,B_n)$ is nonnegative and locally bounded, although the $f_j$ need not be $\theta_j$-psh.
\end{proposition}
\begin{proof}
The potentials $\rho_j$ exist by \cref{lem:smooth-potential-on-ball}. By \cref{lem:metric-laplacian-local-weak-hessian}, all $u_j$ and $f_j$ belong to $W^{2,2}_{\loc}(U)$, so the matrices $A_j$ and $B_j$ are defined almost everywhere, and the positive current $T_j=\theta_j+\ddc u_j=\ddc(\rho_j+u_j)$ has coefficients in $L^2_{\loc}$.

We first show that $A_j$ is positive semidefinite almost everywhere. For each $\xi\in\mathbb C^n$ the distribution $\sum_{a,b}(\rho_j+u_j)_{a\bar b}\xi_a\bar\xi_b$ is positive, because $T_j$ is a positive current; since it is a locally integrable function, it is nonnegative almost everywhere. We apply this to the countably many $\xi$ with rational real and imaginary parts. Outside the union of the corresponding null sets, $\sum_{a,b}(A_j)_{ab}\xi_a\bar\xi_b\geq0$ for all such $\xi$, hence for all $\xi\in\mathbb C^n$, because these $\xi$ are dense and a quadratic form on $\mathbb C^n$ is continuous.

Next we bound the coefficients. By the trace identity \eqref{eq:kahler-trace-star},
\[
 \Lambda T_j=\Lambda\theta_j-\frac{Lu_j}{4\pi},
\]
which is bounded on $W$ by hypothesis. Write $\omega=\frac{\mathrm i}{2\pi}\sum_{a,b}h_{a\bar b}\,\mathrm dz_a\wedge\mathrm d\bar z_b$ on $W$. Since $T_j=\frac{\mathrm i}{2\pi}\sum_{a,b}(A_j)_{ab}\,\mathrm dz_a\wedge\mathrm d\bar z_b$, the function $\Lambda T_j$ is the trace of $A_j$ relative to $h$, as in \eqref{eq:complex-laplacian-trace}. As $\overline W$ is compact, $h\leq CI$ on $W$ for a constant $C$; since $A_j$ is positive semidefinite, it follows that $\operatorname{tr}A_j\leq C\Lambda T_j$ almost everywhere on $W$. For a positive semidefinite Hermitian matrix every diagonal entry is at most the trace, and every entry satisfies $|a_{ab}|^2\leq a_{aa}a_{bb}$. Therefore $A_j$ is bounded on $W$.

The functions $\rho_j+u_j$ are continuous and plurisubharmonic on $W$, and they have bounded complex Hessian in the sense of \cref{sec:bounded-complex-hessian-contact}. The first identity is therefore the first part of \cref{prop:bounded-complex-hessian-density}. The functions $\rho_j+f_j$ are real, continuous and belong to $W^{2,2}_{\loc}(W)$, and adding $\rho_j$ to both $u_j$ and $f_j$ does not change the contact set $\{u_j=f_j\}$. Hence the second part of \cref{prop:bounded-complex-hessian-density} gives the second identity. Moreover \cref{prop:bounded-complex-hessian-contact}, applied to $\rho_j+u_j$ and $\rho_j+f_j$, shows that $A_j=B_j$ almost everywhere on $\{u_j=f_j\}$, hence on $E$. Consequently $\mathbf1_ED(B_1,\ldots,B_n)=\mathbf1_ED(A_1,\ldots,A_n)$ almost everywhere. This density is nonnegative, since by the first identity $D(A_1,\ldots,A_n)\,\mathrm dV_{2n}$ is a positive measure, and it is bounded, since $D$ is a polynomial in the entries of the bounded matrices $A_j$.
\end{proof}
\index{contact set!metric-Laplacian obstacles}

A general obstacle may have unbounded complex-Hessian coefficients even when its metric Laplacian is bounded. The following estimate nevertheless makes all finite products of these coefficients integrable.

\begin{proposition}\label{prop:compact-metric-laplacian-lp}
Let $1<p<\infty$ and let $f\in L^p(X)$ have $Lf\in L^p(X)$ distributionally. Then $f\in W^{2,p}(X)$ and
\begin{equation}\label{eq:compact-metric-laplacian-lp}
 \|f\|_{W^{2,p}(X)}\leq C_{X,\omega,p}
       \bigl(\|Lf\|_{L^p(X)}+\|f\|_{L^p(X)}\bigr).
\end{equation}
In particular a continuous function with bounded distributional metric Laplacian belongs to $W^{2,p}(X)$ for every finite $p\geq1$ and to $C^{1,\alpha}(X)$ for every $0<\alpha<1$.
\end{proposition}
\begin{proof}
We first prove the estimate for smooth real functions. In real coordinates of dimension $d=2n$, write
\[
 L=-\sum_{i,j}a^{ij}(x)\partial_i\partial_j
               +\sum_i b^i(x)\partial_i,
\]
where the symmetric matrix $a$ is positive definite and all coefficients are smooth. On the compact chart supports its eigenvalues have fixed positive upper and lower bounds. For the frozen operator $L_0=-\sum a^{ij}(x_0)\partial_i\partial_j$, a linear change of coordinates reduces the principal form to the Euclidean one. Equation \eqref{eq:euclidean-hessian-lp} gives
\[
 \|D^2v\|_p\leq C\|L_0v\|_p.
\]
The constant is uniform when the eigenvalues of $a(x_0)$ stay in the fixed bounds. On a sufficiently small ball, $\sup|a(x)-a(x_0)|\leq\delta$, and hence
\[
 \|D^2v\|_p\leq C\bigl(\|Lv\|_p+
           \delta\|D^2v\|_p+\|\nabla v\|_p\bigr)
\]
for $v$ compactly supported there. Choose $\delta$ to absorb the Hessian term.

Take a finite smooth partition $\chi_\nu$ subordinate to these balls and apply the estimate to $v_\nu=\chi_\nu f$, extended by zero in its chart. The commutators $[L,\chi_\nu]f$ contain only first derivatives of $f$. Summing gives
\[
 S_2:=\sum_\nu\|D^2v_\nu\|_p
 \leq C\bigl(\|Lf\|_p+\|\nabla f\|_p+\|f\|_p\bigr).
\]
All chart and metric norms here are equivalent on the fixed finite cover. Viewing the localized functions again on $X$, we have $f=\sum_\nu v_\nu$; its gradient norm is at most a fixed multiple of $\sum_\nu\|\nabla v_\nu\|_p$. By \cref{lem:gradient-lp-interpolation}, this is at most $C\varepsilon S_2+C_\varepsilon\|f\|_p$. Absorption gives \eqref{eq:compact-metric-laplacian-lp} for smooth functions, including their first derivatives.

Let $P=L+c$ with the fixed $c>0$ used in the proof of \cref{lem:compact-elliptic-foundations}. For smooth real $b$, that proof gives a smooth real solution $u=P^{-1}b$. Testing the equation by
\[
 \phi_\varepsilon(u)=u(u^2+\varepsilon)^{(p-2)/2},
 \qquad \varepsilon>0,
\]
is legitimate. Its derivative is
\[
 \phi_\varepsilon'(u)=
 ((p-1)u^2+\varepsilon)(u^2+\varepsilon)^{(p-4)/2}\geq0.
\]
Stokes therefore gives $c\int_Xu\phi_\varepsilon(u)\,\mathrm dV\leq\int_Xb\phi_\varepsilon(u)\,\mathrm dV$. Smoothness and compactness permit dominated convergence as $\varepsilon\downarrow0$, also for $1<p<2$, where $|\phi_\varepsilon(u)|\leq|u|^{p-1}$. H\"older gives
\[
 \|P^{-1}b\|_p\leq c^{-1}\|b\|_p.
\]
The smooth estimate already proved, applied with $Lu=b-cu$, now yields $\|P^{-1}b\|_{W^{2,p}}\leq C\|b\|_p$.

For the given real $f$ put $b=Lf+cf\in L^p$, and choose smooth real $b_j\to b$ in $L^p$ by chart cutoffs and mollification. The solutions $u_j=P^{-1}b_j$ are Cauchy in $W^{2,p}$. Completeness follows coefficientwise from completeness of $L^p$, with the weak derivatives identified by testing. Their limit $u$ satisfies $Pu=b$ distributionally. The distributional inverse identities constructed in \cref{lem:compact-elliptic-foundations} give uniqueness for $P$, so $u=f$. This proves the estimate. Apply the real assertion to the real and imaginary parts for complex $f$.

If $f$ is continuous and $Lf$ is bounded, both are in every finite $L^p$ on $X$. The assertion for $p=1$ follows from $W^{2,2}$ and finite volume. Given $0<\alpha<1$, choose $p>d/(1-\alpha)$. Morrey's inequality, \cref{lem:morrey-sobolev-representative}, in each chart gives a $C^{1,1-d/p}$ representative, hence $C^{1,\alpha}$. It equals the original continuous function because the two agree almost everywhere.
\end{proof}
\index{Laplacian!$W^{2,p}$ regularity}

\begin{corollary}\label{cor:metric-laplacian-barrier-determinants}
Let $\theta_1,\ldots,\theta_n$ be closed smooth real $(1,1)$-forms on $X$, and let $f_1,\ldots,f_n$ be continuous real functions with bounded distributional $Lf_j$. Their weak forms have a well-defined wedge
\[
 \bigwedge_{j=1}^n(\theta_j+\ddc f_j)
       =h\,\mathrm dV,\qquad h\in L^q(X)\quad(1\leq q<\infty).
\]
This is a finite signed Radon measure. On a coordinate ball choose smooth real $\rho_j$ with $\theta_j=\ddc\rho_j$; there it is $n!\pi^{-n}D(B_1,\ldots,B_n)\,\mathrm d\lambda$, with $B_j=(\partial_a\bar\partial_b(\rho_j+f_j))$. If continuous $\theta_j$-psh functions $u_j$ also have bounded $Lu_j$, then on $E=\bigcap_j\{u_j=f_j\}$ this signed measure agrees with the positive contact measure in \cref{prop:metric-laplacian-regular-contact-density}.
\end{corollary}
\begin{proof}
The proposition puts all weak second derivatives of $f_j$ in every finite $L^p$. Products of $n$ coefficients are therefore in $L^q$ by H\"older with $p=nq$; when $nq=1$, use $p=2$ and finite volume instead. Expansion of the wedges defines their density. Weak differential forms transform by the smooth coordinate rule, so these polynomial expressions agree on chart overlaps almost everywhere. The normalization is the smooth determinant calculation in \cref{prop:bounded-complex-hessian-density}, applied algebraically to the weak matrices. A real $L^1$ density defines a finite signed Radon measure by its positive and negative parts. The contact identity is precisely \cref{prop:metric-laplacian-regular-contact-density}.
\end{proof}
\index{Monge--Amp\`ere measure!weak barrier determinant}

